\documentclass[a4paper,12pt]{article}
\usepackage[T2A]{fontenc}
\usepackage[utf8]{inputenc}
\usepackage[russian]{babel}
\usepackage{amsmath,amssymb,mathtools}
\usepackage{icomma}
\usepackage[left=22mm,right=22mm,top=22mm,bottom=22mm]{geometry}
\usepackage{microtype}
\usepackage{tikz}
\usetikzlibrary{calc,arrows.meta,positioning,shapes.geometric}
\usepackage{booktabs,tabularx}
\usepackage{enumitem}
\usepackage{parskip}
\usepackage{fancyhdr}
\usepackage{setspace}
\usepackage{xcolor}
\usepackage[hidelinks]{hyperref}
\definecolor{accent}{HTML}{008080}
\definecolor{darkaccent}{HTML}{004D4D}
\definecolor{textgray}{HTML}{333333}
\definecolor{quiet}{HTML}{657577}
\onehalfspacing
\setlength{\parindent}{0pt}
\setlength{\parskip}{5pt}
\setlist[itemize]{nosep,leftmargin=1.65em,topsep=4pt}
\setlist[enumerate]{leftmargin=1.85em,itemsep=7pt,topsep=8pt}
\renewcommand{\arraystretch}{1.22}
\pagestyle{fancy}
\fancyhf{}
\fancyhead[L]{\footnotesize\color{quiet} Неравенства со степенями и логарифмами}
\fancyhead[R]{\footnotesize\color{quiet} Тимофей Васильченко}
\fancyfoot[C]{\small\thepage}
\renewcommand{\headrulewidth}{0.25pt}
\setlength{\headheight}{15pt}
\newcommand{\lead}[1]{\textcolor{darkaccent}{\textbf{#1}}}
\newcommand{\important}[1]{\textcolor{accent}{#1}}
\newcommand{\ts}{\textsuperscript{\small\textcolor{quiet}{*}}}
\newcommand{\work}[1]{\par\vspace{3pt}\textcolor{darkaccent}{\large\bfseries #1}\par\vspace{3pt}}
\hypersetup{pdftitle={Чек-лист: показательные и логарифмические неравенства},pdfauthor={Лёвин Артём Александрович}}
\begin{document}
\hypersetup{pageanchor=false}
\begin{titlepage}
\thispagestyle{empty}
\centering
\vspace*{22mm}
{\large\color{darkaccent} МАТЕМАТИКА}\par
\vspace{17mm}
{\LARGE Чек-лист}\par
\vspace{14mm}
{\LARGE\color{accent} Дробные неравенства\\[5pt] со степенями и логарифмами}\par
\vspace{15mm}
{\large Подготовка к ЕГЭ по профильной математике}\par
\vfill
{\large Тимофей Васильченко}\par
\vspace{12mm}
{6 ноября 2026 г.}\par
\vspace{15mm}
{\small Лёвин Артём Александрович\\[2pt] эксклюзивно для Тимофея Васильченко}\par
\vspace{17mm}
\begin{tikzpicture}[x=1cm,y=1cm]
\draw[accent, line width=.8pt] (0,0) .. controls (1.2,.55) and (2.0,-.52) .. (3.2,0)
 .. controls (4.4,.52) and (5.2,-.55) .. (6.4,0)
 .. controls (7.5,.49) and (8.25,-.35) .. (9,0);
\end{tikzpicture}
\end{titlepage}
\setcounter{page}{1}
\hypersetup{pageanchor=true}

\section*{1. Что именно решаем}
Неравенство содержит дробь: в числителе --- разность показательных выражений, в знаменателе --- квадрат логарифма. Цель --- определить, где дробь имеет нужный знак, \emph{не нарушая область допустимых значений (ОДЗ)}.
\[
\frac{6^{5-2x}-0{,}5^{5-2x}}{1-\log_{0{,}04}^{2}(2x-5)}\ge 0.
\tag{1}\label{eq:main}
\]
Здесь $\log_a^2 t=(\log_a t)^2$; знак квадрата относится ко всему значению логарифма.

\section*{2. ОДЗ --- раньше преобразований}
Для $\log_{0{,}04}(2x-5)$ аргумент должен быть строго положительным:
\[
2x-5>0\quad\Longrightarrow\quad \boxed{x>\frac52}.
\]
Кроме этого, знаменатель дроби \emph{не должен обращаться в нуль}. Его нули найдём ниже и исключим. Условия $a>0$, $a\ne1$ для основания логарифма выполнены: $0{,}04=1/25$.

\section*{3. Где числитель обращается в нуль}
Приравниваем числитель к нулю:
\begin{align*}
6^{5-2x}-0{,}5^{5-2x}&=0,\\
6^{5-2x}&=0{,}5^{5-2x},\\
\left(\frac6{0{,}5}\right)^{5-2x}&=1,\\
12^{5-2x}&=12^0,\\
5-2x&=0\quad\Longrightarrow\quad x=\frac52.
\end{align*}
Деление законно: $0{,}5^{5-2x}>0$ при любом $x$. Но точка $x=5/2$ \lead{не входит в ОДЗ}. Следовательно, на допустимой области числитель не равен нулю. Более того, при $x>5/2$ имеем $5-2x<0$ и
\[
6^{5-2x}<0{,}5^{5-2x},\qquad
\boxed{6^{5-2x}-0{,}5^{5-2x}<0}.
\]
Это полезный способ сразу установить знак числителя на всей ОДЗ.

\section*{4. Нули знаменателя}
\begin{align*}
1-\log_{0{,}04}^{2}(2x-5)&=0,\\
\log_{1/25}^{2}(2x-5)&=1,\\
\log_{1/25}(2x-5)&=1\quad\text{или}\quad
\log_{1/25}(2x-5)=-1.
\end{align*}
В первом случае $2x-5=1/25$, во втором $2x-5=25$. Получаем:
\[
\boxed{x=\frac{63}{25}=2{,}52}\qquad\text{и}\qquad\boxed{x=15}.
\]
\lead{Обе точки исключаются}, поскольку при них знаменатель равен нулю. Полная допустимая область:
\[
x>\frac52,\qquad x\ne\frac{63}{25},\qquad x\ne15.
\]

\section*{5. Метод интервалов: проверяем знаки}
С учётом ОДЗ остаются три промежутка:
\[
\left(\frac52,\frac{63}{25}\right),\qquad
\left(\frac{63}{25},15\right),\qquad(15,+\infty).
\]
Удобно начать со средней области и взять $x=10$.
\begin{itemize}
\item Числитель отрицателен: $6^{-15}-0{,}5^{-15}=\frac1{6^{15}}-2^{15}<0$.
\item Поскольку $0{,}04=25^{-1}$, то $\log_{0{,}04}15=-\log_{25}15$.
      Из $1<15<25$ следует $0<\log_{25}15<1$.
      Поэтому $1-(\log_{25}15)^2>0$.
\end{itemize}
Итак, при $x=10$ дробь отрицательна. На левом крайнем промежутке $0<2x-5<1/25$, поэтому $\log_{1/25}(2x-5)>1$ и знаменатель отрицателен. На правом $2x-5>25$, поэтому логарифм меньше $-1$ и знаменатель также отрицателен. Числитель везде отрицателен, следовательно, на обоих крайних промежутках дробь положительна.

\begin{center}
\begin{tikzpicture}[x=1cm,y=1cm,>=Latex]
  \draw[->,very thick,accent] (0,0)--(13.0,0) node[right] {$x$};
  \foreach \x/\lab in {1.0/{$2{,}5$},5.0/{$2{,}52$},9.8/{$15$}}
    {\draw[white,fill=white] (\x,0) circle (2.6pt); \draw[accent,thick,fill=white] (\x,0) circle (2.6pt); \node[below=6pt] at (\x,0) {\lab};}
  \node[above=7pt,quiet] at (3.0,0) {$+$};
  \node[above=7pt,quiet] at (7.4,0) {$-$};
  \node[above=7pt,quiet] at (11.2,0) {$+$};
  \node[below=6pt,quiet] at (.35,0) {вне ОДЗ};
\end{tikzpicture}
\end{center}
Все граничные точки выколоты. Поскольку нуль числителя лежит вне ОДЗ, знак $\ge0$ не даёт включённых точек.
\[
\boxed{\displaystyle x\in\left(\frac52;\frac{63}{25}\right)\cup(15;+\infty)}.
\]
\lead{Важное правило.} Знаки следует устанавливать на допустимых промежутках. Автоматически чередовать их можно только после обоснования; сама по себе расстановка точек на оси этого не гарантирует.

\section*{6. Альтернативный путь: убрать степени и логарифмы из анализа знака}
Это второй метод, рассмотренный на занятии. Он особенно полезен, когда прямые подстановки дают громоздкие степени и логарифмы.

\work{Два правила сравнения знаков}
Для $a>0$, $a\ne1$ и действительных $u,v$:
\[
\operatorname{sgn}(a^u-a^v)=\operatorname{sgn}\bigl((a-1)(u-v)\bigr).
\tag{2}
\]
Для $b>0$, $c>0$ при том же условии на $a$:
\[
\operatorname{sgn}(\log_a b-\log_a c)
 =\operatorname{sgn}\bigl((a-1)(b-c)\bigr).
\tag{3}
\]
При $a>1$ обе функции возрастают; при $0<a<1$ убывают. Поэтому множитель $a-1$ правильно учитывает направление сравнения.

\work{Разложение исходного выражения}
Обозначим $t=2x-5$; на ОДЗ $t>0$. Вынесем положительный множитель из числителя и разложим знаменатель как разность квадратов:
\[
\frac{0{,}5^{5-2x}\bigl(12^{5-2x}-1\bigr)}
{\bigl(1-\log_{1/25}t\bigr)\bigl(1+\log_{1/25}t\bigr)}\ge0.
\]
Первый множитель числителя строго положителен. По правилу~(2) знак скобки $12^{5-2x}-12^0$ совпадает со знаком $5-2x$.

\work{Логарифмы заменяем сравнениями}
Из $\log_{1/25}t=-\log_{25}t$ и $1=\log_{25}25$ получаем:
\begin{align*}
1-\log_{1/25}t&=\log_{25}25+\log_{25}t
=\log_{25}25-\log_{25}\!\frac1t,\\
1+\log_{1/25}t&=\log_{25}25-\log_{25}t.
\end{align*}
Так как основание $25>1$, знаки этих множителей совпадают соответственно со знаками $25-1/t$ и $25-t$.

Поэтому для \emph{допустимых} $x$ исходная дробь имеет тот же знак, что и
\[
\frac{5-2x}{\displaystyle\left(25-\frac1{2x-5}\right)
\bigl(25-(2x-5)\bigr)}.
\tag{4}
\]
Сохраняем исключения $x=63/25$ и $x=15$; преобразование касается знаков, а не равенства значений двух дробей.

\work{Сведение к дробно-рациональному неравенству}
Поскольку $5-2x=-t$ и $t>0$, выражение~(4) можно переписать:
\[
\frac{-t^2}{(25t-1)(25-t)}\ge0,\qquad t>0,
\qquad t\ne\frac1{25},\ 25.
\]
Числитель отрицателен. Значит, знаменатель должен быть отрицателен:
\[
(25t-1)(25-t)<0
\quad\Longleftrightarrow\quad
0<t<\frac1{25}\ \text{или}\ t>25.
\]
Возвращаясь к $t=2x-5$, снова получаем
\[
\boxed{\displaystyle x\in\left(\frac52;\frac{63}{25}\right)\cup(15;+\infty)}.
\]
Оба метода дают один ответ. Первый особенно удобен, если легко подобрать пробные точки; второй --- если удалось разложить выражения на множители и заменить их эквивалентными \emph{по знаку} простыми выражениями.

\section*{7. Точные вычисления и запись}
В упражнении на занятии возникла граница
\[
4-5x=\frac18.
\]
Её следует находить без преждевременного округления:
\[
5x=4-\frac18=\frac{31}{8},\qquad
\boxed{x=\frac{31}{40}=0{,}775}.
\]
Если одновременно имеется условие $4-5x>0$, то ОДЗ требует $x<4/5=0{,}8$. Числа $0{,}775$ и $0{,}8$ различны и должны быть верно расположены на числовой оси.

\lead{Проверяйте запись:}
\begin{itemize}
\item В логарифме различайте $\bigl(\log_a t\bigr)^2$ и $\log_a(t^2)$.
\item Квадрат отрицательного числа записывайте со скобками: $(-1)^2=1$.
\item Для ограничения $4-5x>0$ после деления на $-5$ знак меняется: $x<4/5$.
\item Приближённые десятичные значения не заменяют точные границы ответа.
\item Нули знаменателя выколоты даже в нестрогом неравенстве.
\end{itemize}

\clearpage
\thispagestyle{plain}
\vspace*{6mm}
\begin{center}
{\Large\color{darkaccent}\bfseries Алгоритм: метод интервалов для дроби}\par
\vspace{11mm}
\begin{tikzpicture}[node distance=10mm and 16mm,>=Latex,
    every node/.style={align=center,font=\small},
    actbox/.style={draw=accent, rounded corners=3pt, line width=.65pt, text width=9.6cm, inner sep=7pt},
    oval/.style={draw=darkaccent, ellipse, line width=.8pt, text width=4.2cm, inner sep=4pt},
    decision/.style={draw=accent,diamond,aspect=2.8,text width=4.6cm,inner sep=3pt},
    arrow/.style={-{Latex[length=2.5mm]},thick,draw=darkaccent}]
  \node[oval] (start) {Начать с дробного неравенства};
  \node[actbox,below=of start] (dom) {Найти ОДЗ: аргументы логарифмов $>0$; знаменатель $\ne0$};
  \node[actbox,below=of dom] (roots) {Найти нули числителя и знаменателя, определить допустимые точки};
  \node[actbox,below=of roots] (line) {Отметить на оси границы и выделить промежутки внутри ОДЗ};
  \node[decision,below=12mm of line] (cond) {Есть допустимые промежутки?};
  \node[actbox,below=12mm of cond] (signs) {Найти знак дроби на каждом промежутке: пробной точкой или анализом множителей};
  \node[actbox,below=of signs] (select) {Оставить промежутки нужного знака; допустимые нули числителя включить при $\ge,\le$};
  \node[oval,below=of select] (finish) {Записать ответ};
  \node[actbox,right=19mm of cond,text width=2.55cm,inner sep=5pt] (empty) {Ответ: $\emptyset$};
  \draw[arrow] (start)--(dom);
  \draw[arrow] (dom)--(roots);
  \draw[arrow] (roots)--(line);
  \draw[arrow] (line)--(cond);
  \draw[arrow] (cond)--node[right,fill=white,inner sep=1pt]{да} (signs);
  \draw[arrow] (cond)--node[above,fill=white,inner sep=1pt]{нет} (empty);
  \draw[arrow] (signs)--(select);
  \draw[arrow] (select)--(finish);
\end{tikzpicture}
\end{center}
\clearpage

\section*{Тренировочный блок --- Путь А}
\textcolor{quiet}{10 задач. От точных вычислений и поиска ОДЗ --- к самостоятельному решению смешанных неравенств.}\par

\begin{enumerate}[itemsep=17pt]
\item \textbf{Точные границы.} Вычислите без округления, запишите обыкновенной и десятичной дробями:
\[
\frac{4-\frac18}{5}.
\]

\item \textbf{ОДЗ.} Найдите область определения выражения
\[
\log_{1/25}(2x-5).
\]

\item \textbf{Нуль числителя.} Решите уравнение
\[
8^{4-2x}-0{,}5^{4-2x}=0.
\]

\item \textbf{Нули знаменателя.} Решите уравнение
\[
1-\log_{1/16}^{2}(3x-2)=0.
\]

\item \textbf{Знак без вычисления значения.} Определите знак дроби при $x=0$:
\[
\frac{3^{2x-4}-0{,}5^{2x-4}}{1-\log_{1/9}^{2}(x+4)}.
\]

\end{enumerate}
\clearpage
\section*{Тренировочный блок --- продолжение}
\textcolor{quiet}{Задачи 6--10. Работайте в точных дробях, отмечайте ОДЗ и выколотые точки.}\par
\begin{enumerate}[start=6,itemsep=26pt]
\item \textbf{Квадрат логарифма.} Решите неравенство
\[
1-\log_{1/9}^{2}(2x+1)>0.
\]

\item \textbf{ОДЗ и два нуля знаменателя.} Решите неравенство
\[
\frac{7^{4-3x}-0{,}5^{4-3x}}
{1-\log_{1/16}^{2}(3x-4)}\ge0.
\]

\item \textbf{Знак смешанной дроби.} Решите неравенство
\[
\frac{2^{1-2x}-0{,}25^{1-2x}}
{1-\log_{1/25}^{2}(x-1)}<0.
\]

\item \textbf{Три граничные точки.} Решите неравенство
\[
\frac{4^{3-x}-0{,}5^{3-x}}
{1-\log_{1/36}^{2}(5-x)}>0.
\]

\item \textbf{Нуль числителя внутри ОДЗ.} Решите неравенство
\[
\frac{2^{x-1}-0{,}25^{x-1}}
{1-\log_{1/4}^{2}(4-2x)}\le0.
\]
\end{enumerate}

\clearpage
\section*{Ответы к тренировочному блоку}
\begin{center}
\begin{tabularx}{\linewidth}{@{}>{\bfseries}p{.7cm}X@{}}
\toprule
№ & Ответ \\
\midrule
1 & $\dfrac{31}{40}=0{,}775$. \\
2 & $x>\dfrac52$. \\
3 & $x=2$. \\
4 & $x=\dfrac{11}{16}$ или $x=6$. \\
5 & Дробь отрицательна. \\
6 & $\displaystyle x\in\left(-\dfrac49;4\right)$. \\
7 & $\displaystyle x\in\left(\dfrac43;\dfrac{65}{48}\right)\cup\left(\dfrac{20}{3};+\infty\right)$. \\
8 & $\displaystyle x\in\left(\dfrac{26}{25};26\right)$. \\
9 & $\displaystyle x\in(-31;3)\cup\left(\dfrac{179}{36};5\right)$. \\
10 & $\displaystyle x\in(0;1]\cup\left(\dfrac{15}{8};2\right)$. \\
\bottomrule
\end{tabularx}
\end{center}

\vspace{16mm}
\section*{Итоговая памятка}
\begin{tabularx}{\linewidth}{@{}p{.36\linewidth}X@{}}
\toprule
Сигнал & Ваше действие \\
\midrule
Логарифм в знаменателе & Сначала ОДЗ и отдельно нули знаменателя. \\
Одинаковые показатели & Разделите на положительную степень; объедините основания. \\
$1-\log_a^2 t$ & Разность квадратов; границы при $\log_a t=\pm1$. \\
Сложный знак & Сравнивайте с нулём; точное значение обычно излишне. \\
Отрицательная степень & $a^{-n}=1/a^n$; проверяйте положительность основания. \\
Нестрогий знак & Включайте только допустимые нули числителя. \\
Десятичные границы & Сохраняйте точные дроби до записи ответа. \\
\bottomrule
\end{tabularx}
\vfill
\begin{center}\footnotesize\color{quiet}Лёвин Артём Александрович --- эксклюзивно для Тимофея Васильченко\end{center}
\end{document}
